\chapter{Allocation and reference counting}
\label{ch:extents}

Until this point blocks came from a bump cursor: a high-water mark in
\code{struct bitter\_root}, advanced on every allocation and never moved back.
Nothing was ever reclaimed, and a leaf that emptied held its block forever.
That was fine while no workload deleted enough to notice.

This chapter is what replaces it. An \emph{extent tree} records which ranges of
the device are in use and how many things point at each; an allocator hands out
addresses from what is left over; and a reference count decides when a range
becomes available again. Refcounting is also what will make snapshots work, and
what makes them safe to delete --- a snapshot is nothing but a second reference
to a subtree.

It is the one place in the filesystem where an off-by-one silently corrupts
data. A range freed while still referenced is handed out again and overwritten,
and the damage is invisible until something reads what used to be there.


\section{The extent tree}

Extent items are keyed \keytriple{bytenr}{EXTENT\_ITEM}{length}: the range's
start address is the objectid, its length is the offset. The payload is a
reference count and a flags word, and nothing else --- the key already says
which range is being described.

Two properties follow from that choice, and both are load-bearing.

\paragraph{Absence means free.} An extent has an item if and only if something
references it. The last reference does not write a zero; it deletes the item.
Free space is therefore not stored anywhere --- it is the complement of what is
stored, and the filesystem never has to keep two representations of the same
fact in step.

\paragraph{The tree is in device order.} Keys sort by objectid first, and here
the objectid is an address. So walking the extent tree from its first item to
its last visits every allocated range in ascending order, and the gaps between
consecutive items are exactly the free ranges, already sorted. That turns the
mount-time scan into a single linear pass with no sorting and no index.


\section{The free-space map}

\code{struct bitter\_fs\_info} holds an array of
\code{struct bitter\_free\_extent}, each a \code{\{start, length\}} pair in
device order. It is built once, at mount, by walking the extent tree and
recording the gaps.

It is a \emph{cache}. The extent tree is the durable truth; the map is a
restatement of it that is faster to allocate from, thrown away at unmount and
rebuilt at the next one. Nothing about it reaches the disk, which is why
\code{bitter\_free\_extent} lives in \code{core/fs.h} and not in the format
header: no packing, no little-endian accessors, no layout assertions. Getting
its representation wrong costs a recompile, not a migration.

\paragraph{Every allocation is one block.} That collapses the classic allocator
policy question. First-fit, best-fit and worst-fit differ only when requests
vary in size, so with a uniform \bytes{4096} they produce identical
fragmentation and the simplest wins: take the first range, advance its start,
shrink its length, and drop the entry when it reaches zero.

\paragraph{Taking from the front is the right end.} It allocates low addresses
first, keeping the used region compact and preferentially consuming the
single-block holes that deletes leave behind --- each of which \emph{shortens}
the map when it is exhausted. Taking from the back would be cheaper to remove
but would grow the filesystem outward and leave the holes forever.

\paragraph{An empty map is the out-of-space condition.} Not an arithmetic
comparison against the device size, just \code{free\_count == 0}. Which is both
cheaper and more honest than the bump allocator's
\code{next\_free + BLOCK > total\_bytes}, an expression that had been copied
into nine call sites so each could distinguish \code{ENOSPC} from
\code{ENOMEM}.


\section{Describing the whole device}

The scan derives free space by subtraction, so it is only as good as what the
extent tree describes. If a block is in use and has no item, the scan reports
it as free and the allocator hands it out.

mkfs therefore records what it writes: an item covering the reserved prefix and
the superblock, and one for each of the two tree blocks it lays down. With
those present the extent tree accounts for every byte of the device, the scan
needs no hardcoded constants, and \code{core/} learns the layout by reading it
rather than by being told.

\paragraph{What happened before that was true.} A freshly formatted image
reported one free range covering the entire device from address 0, and the
allocator began handing out the reserved prefix, the superblock, and the two
trees. Nothing failed at allocation time --- every write succeeded. It surfaced
as a checker rejecting a block on the 118th insert, several calls after the
allocation that had overwritten it. The symptom was a corrupt tree; the cause
was a correct allocator obeying a map that was telling the truth about a
filesystem which had not described itself.

The generalisation is worth keeping: a wrong derivation does not produce an
error, it produces a plausible answer.


\section{The recursion}

Recording an allocation means adding a reference to the extent tree. Adding a
reference means modifying the extent tree, which copy-on-writes its blocks,
which allocates --- which needs another reference recorded. The cycle does not
obviously terminate, and each level of it mutates the tree the level above is
walking.

\paragraph{Delayed references break it.} An allocation does not touch the
extent tree at all. It appends a record --- address, length, $\pm 1$ --- to an
array on the open transaction, and the whole set is applied at commit, at a
point where allocating is no longer reentrant.

That is safe because the in-memory free map has \emph{already} been updated.
The block leaves the free list the instant it is handed out, so nothing else
can be given the same address before the set is drained. Only the durable
record lags, and only until commit.

\paragraph{The signature enforces it.} \fnname{trans\_add\_delayed\_ref} takes
no environment. It cannot read a block, cannot descend a tree, and cannot
allocate, so the recursion is impossible rather than merely avoided. The same
device as passing a null transaction to \fnname{btree\_search} to prove a
lookup cannot write.

\paragraph{Merging.} Recording looks for an existing entry with the same
address and length and adds to its delta, and an entry that reaches zero is
removed. So a block allocated and then freed inside one transaction never
reaches the extent tree at all --- which is the common case for a split that is
later abandoned. Without merging the tree would be copy-on-written twice to
create an item and immediately delete it.

The match must be exact. \code{[0x1000, 0x2000)} and \code{[0x0, 0x11000)} are
different items even though one range contains the other, and treating
containment as a match would apply a delta to an item that does not exist.

\paragraph{Draining converges.} Applying an entry copy-on-writes the extent
tree, which allocates, which appends new entries to the set being drained. So
the drain loops until the set is empty rather than iterating a saved count, and
each entry is removed \emph{before} it is applied, so the array can be appended
to safely while one of its entries is in flight.

It terminates because copy-on-write is idempotent within a transaction: once a
block carries the current generation it is modified in place and allocates
nothing. Each pass therefore touches fewer blocks than the last, and eventually
one touches none. That is an argument rather than a proof, so the loop carries
a bound --- a convergence claim that turns out to be wrong should fail where it
can be read, not hang inside a commit.


\section{Commit ordering}

The transaction commit gains two steps ahead of the ones it already had:

\begin{enumerate}
  \item drain the delayed-ref set into the extent tree;
  \item record the extent tree's moved root in the root tree;
  \item repeat until the set is empty \emph{and} the extent root stops moving;
  \item \code{writeback\_all}, then \code{flush\_device};
  \item overwrite the superblock;
  \item \code{flush\_device} again.
\end{enumerate}

Steps 1--3 must precede step 4 because draining dirties extent-tree blocks, and
must precede step 5 because the extent tree's root moves while it runs. Step 2
copy-on-writes the \emph{root} tree, which allocates, which refills the set ---
which is why 1 and 2 are a loop rather than a sequence.

Steps 4--6 are unchanged from the transaction chapter: the first flush is the
ordering barrier that stops the superblock reaching the disk before the blocks
it points at, and the single superblock overwrite is the moment the transaction
becomes real.


\section{What is not reclaimed}

A block freed during a transaction is still referenced by the \emph{committed}
filesystem --- the one the superblock points at, the one that survives a crash
at this instant. Returning it to the free map would let the allocator hand it
out and overwrite it, and a crash before the superblock write would then leave
an intact old tree reading a block containing something else. A filesystem that
was consistent would be destroyed by a transaction that never completed.

btrfs solves this by \emph{pinning}: extents freed during a transaction go to a
separate set and join the free map only after the commit completes. bitterfs
does something simpler and strictly weaker --- it never returns space to the
map at all during a mount. The map only shrinks; freed space becomes available
at the next mount, when the scan rebuilds it from an extent tree that no longer
mentions those ranges.

That makes the hazard structurally impossible rather than carefully avoided,
which is the right trade while the mechanism that would get it right does not
exist. The cost is real and bounded: a long-running mount that deletes heavily
can report \code{ENOSPC} with most of the device unreferenced.

\decide{The specific temptation to resist is rebuilding the map at the end of
the drain. The map it produced would be perfectly consistent with the extent
tree --- and would resurrect every block freed this transaction, into a
transaction that is still allocating.}


\section{Checking it}

A reference count that drifts is invisible until a block is freed while still
referenced, and by then the evidence is gone. That is why
\fnname{bitter-fsck} belongs to this phase rather than a later one.

Its central question is a comparison: what the trees actually point at, versus
what the extent tree claims. It walks every block of every tree with an
explicit worklist --- recording one observation per \emph{pointer}, so the
count falls out of how many times an address appears --- then sorts the
observations and merges them against the extent tree's items, which are already
in address order.

Three findings fall out of that single pass:

\begin{itemize}
  \item a count that disagrees with the number of pointers found;
  \item an item nothing points at, which is a leak: space lost until the
        filesystem is rebuilt;
  \item a block pointed at with no item at all, which is the dangerous
        direction --- the allocator believes that range is free.
\end{itemize}

Walking every block also makes the per-block checks available for the first
time on the whole image rather than on one search path: \field{level} against
what the parent claimed, \field{generation} against the duplicate in the
\code{key\_ptr}, and \field{owner} against the tree being walked. That last one
is the only reader the field has ever had, and it detects a block cross-linked
into two trees --- intact, correctly addressed, and belonging to somebody else,
which the checksum and the self-address check both miss by design.

\paragraph{fsck is read-only,} and not merely out of caution. A tool run to
find damage must not be able to overwrite the evidence. It reports and counts;
it repairs nothing. A repair that is wrong is worse than no repair, and the way
to earn confidence in one is to have a detector that is already trusted.
